\chapter{Risks and Mitigation Plan}

\begin{table}[H]
\centering
\footnotesize
\setlength{\tabcolsep}{4pt}
\renewcommand{\arraystretch}{1.2}

\begin{tabular}{|c|p{3.2cm}|c|p{3.2cm}|p{3.2cm}|p{3.2cm}|}
\hline
\textbf{Rank} & \textbf{Risk} & \textbf{Impact} & \textbf{Mitigation} & \textbf{Response} & \textbf{Contingency} \\
\hline

1 
& Unclear research problem or gap leads to poorly defined research questions 
& High 
& Conduct a comprehensive literature review on misinformation detection models. 
& Reassess the research direction and review the main findings from the literature. 
& Clarify the research objective, narrow the scope, and sharpen the research questions. \\
\hline

2 
& Complex methodology design without a benchmark from existing detection models 
& High 
& Align the design with the proposal scope and benchmark it against prior research. 
& Simplify the framework, model parameters, or components. 
& Adopt baseline models from prior research with modified parameters if needed. \\
\hline

3 
& Difficulty selecting datasets for social-context-aware multimodal analysis 
& High 
& Identify validated datasets used in prior research. 
& Evaluate alternative datasets that may be compatible with the proposed model. 
& Choose alternative datasets aligned with the proposed features and parameters, or adjust the model to fit the selected datasets. \\
\hline

4 
& Time constraints lead to an incomplete research proposal 
& High 
& Set weekly milestones and track progress each week. 
& Reallocate tasks to team members with available capacity. 
& Prioritise core sections and work to a strict day-by-day timeline. \\
\hline

5 
& Inconsistent writing quality across team members 
& Medium 
& Assign clear section ownership to each member. 
& Conduct peer review and standardise writing style. 
& Complete final integration and editing through a designated member. \\
\hline

6 
& Uneven workload distribution among team members 
& Medium 
& Allocate tasks clearly and monitor progress weekly. 
& Reassign tasks during weekly meetings or sprint reviews. 
& Adjust future workload distribution based on progress and capacity. \\
\hline

\end{tabular}

\caption{Risk Management Plan}
\label{tab:risk_management_plan}
\end{table}



\chapter{Generative AI Usage}
 
% ─────────────────────────────────────────────────────────────
% DECLARATION
% Replace the block below with the appropriate declaration from
% https://aiattribution.github.io/ once you know which tools
% were used. Two examples are provided — keep only the one
% that applies, or replace with your generated statement.
% ─────────────────────────────────────────────────────────────
 
\section*{Declaration}
 
% Example A — No AI used:
% \textit{This work was entirely human-created, without the use of AI.} AIA No AI v1.0
 
% Example B — AI used (update tool name(s) as appropriate):
% \textit{This work was primarily human-created. AI was used to make stylistic edits,
% such as changes to structure, wording, and clarity. AI was used to make content edits,
% such as changes to scope, information, and ideas. AI was prompted for its contributions,
% or AI assistance was enabled. AI-generated content was reviewed and approved.
% The following model(s) or application(s) were used: [Tool Name].} AIA No AI v1.0
 
\noindent \textbf{[REPLACE THIS LINE WITH YOUR GENERATED DECLARATION FROM \url{https://aiattribution.github.io/}]}
 
% ─────────────────────────────────────────────────────────────
% AI TOOL REPORT TABLE(S)
% Duplicate this entire section for each additional tool used.
% Delete this section entirely if no AI tools were used.
% ─────────────────────────────────────────────────────────────
 

\section*{Tool Report: [Tool Name, e.g. ChatGPT / RMIT Val / Copilot]}
 
\begin{table}[H]
\renewcommand{\arraystretch}{1.4}
\begin{tabular}{|p{0.35\linewidth}|p{0.58\linewidth}|}
    \hline
    \textbf{Field} & \textbf{Details} \\
    \hline
 
    1.\ GenAI tool used
    & [e.g. ChatGPT 4o, RMIT Val, GitHub Copilot] \\
    \hline
 
    2.\ Purpose for this assignment
    & [e.g. Brainstorming research questions, drafting outline,
      rewriting for clarity, code generation] \\
    \hline
 
    3.\ Prompts or instructions given
    & [Include key prompts used. If conversation was lengthy, include
      all your prompts in full and use {[}\ldots{]} to truncate AI
      responses only.] \\
    \hline
 
    4.\ Number of prompts given
    & [e.g. 7] \\
    \hline
 
    5.\ Summary of GenAI output
    & [Describe what the AI produced. Include brief excerpts where
      relevant.] \\
    \hline
 
    6.\ Verification of GenAI output
    & [Explain how you checked the AI's output for factual accuracy,
      source validity, and logical soundness. e.g. ``All cited sources
      were manually located and verified via Google Scholar.''] \\
    \hline
 
    7.\ Usage or modification of output
    & [Explain how you used, adapted, or discarded the AI-generated
      content in your final submission.] \\
    \hline
 
    8.\ Reflection
    & [What worked well? What did not? How did it support your
      learning? What value did it add?] \\
    \hline
 
\end{tabular}
\caption{Generative AI usage report — [Tool Name]}
\end{table}
 
% ─────────────────────────────────────────────────────────────
% Add another \section* and table block here for each
% additional tool used. e.g.:
%
% \section*{Tool Report: GitHub Copilot}
% \begin{table}[h!] ... \end{table}
% ─────────────────────────────────────────────────────────────